\scanpage{27}
que l'opération de $\pi$ sur $\tau$ a été prolongée en une opération de $\mathcal E$ sur $\tau$. Mais ceci ne donne pas d'interprétation \unsure{symp.} de la $T$-structure\ldots
\figure{\begin{tikzpicture}\draw (0,0)--(1.5,0);\end{tikzpicture}}

\textcircled{9} Les schémas et Topos fondamentaux dans la théorie.

\textcircled{A} \deleted{Topos \unsure{Etendus|Etendue}} \inserted{Topos} modulaires \inserted{$(\mathcal T_{0,\nu})_{\overline{\mathbf Q}}$} pour les courbes de genre $0$ avec $\nu$ points marqués, et famille \deleted{des} modulaire \deleted{\unsure{par}} $\nu$-ponctuée sur $M_{0\nu}$\ldots

Le Topos $(\mathcal T_{0,\nu})_{\mathbf Q}$ pour $\nu=0,1,2$ \uline{n'est} pas un étendu --- car le schéma d'Isom \unsure{d'un|d'une} $X/S$ fibré projectif de rang $1$, avec $\nu$ sections marquées (distinctes sur chaque fibre) est lisse de dim rel. $3-\nu\deleted{\unsure{\geq|=}}>0$ (si $\deleted{\unsure{\nu}}0\leq\nu\leq2$).
\deleted{Néanmoins} \inserted{En fait} ce sont les topos
\[
B_{\unsure{\operatorname{GP}(1)|\operatorname{PGL}(1)}},\quad B_{\operatorname{Aff}(1)},\quad B_{G_m},\quad\inserted{B_e=e}
\]
dont les
\mspage{26}
